\chapter{Kesimpulan dan Saran}

\section{Kesimpulan}

Pada tugas besar ini telah diimplementasikan enam algoritma \textit{local search}, yaitu empat varian \textit{Hill-climbing} (\textit{steepest ascent}, \textit{sideways move}, \textit{random restart}, dan \textit{stochastic}), \textit{Simulated Annealing}, dan \textit{Genetic Algorithm}, untuk menyelesaikan permasalahan pemuatan paket ke dalam truk yang bertujuan memaksimalkan total nilai paket yang dimuat. \textit{State} direpresentasikan sebagai posisi dan orientasi setiap paket, sedangkan seluruh \textit{constraint} (batas dinding, tumpang tindih, kapasitas berat, paket melayang, dan paket rapuh) diperlakukan sebagai \textit{hard constraint}. Karena \textit{state} awal dan seluruh \textit{neighbor} yang dihasilkan selalu diperiksa sebelum digunakan, fungsi objektif cukup berupa jumlah nilai paket yang dimuat tanpa suku penalti. Keenam algoritma memakai komponen yang sama (representasi \textit{state}, fungsi objektif, operator \textit{neighbor}, dan format hasil), sehingga perbandingan antaralgoritma dapat dilakukan secara adil dan seluruh data yang diminta spesifikasi dapat dicatat dengan cara yang seragam.

Ketiga kelompok algoritma menunjukkan karakteristik yang berbeda. \textit{Hill-climbing} hanya berpindah ke \textit{neighbor} yang lebih baik sehingga pencarian berhenti pada \textit{local optimum}. Varian \textit{sideways move} dan \textit{random restart} berupaya mengurangi masalah tersebut, masing-masing dengan melewati daerah datar dan memulai ulang dari \textit{state} awal yang baru, tetapi tidak menjamin tercapainya \textit{global optimum}. \textit{Simulated Annealing} dapat menerima \textit{neighbor} yang lebih buruk dengan probabilitas $e^{\Delta E/T}$ yang mengecil seiring turunnya suhu, sehingga mampu keluar dari \textit{local optimum} pada tahap awal dan bersifat semakin selektif pada tahap akhir. \textit{Genetic Algorithm} menjaga sekumpulan solusi sekaligus dan menggabungkannya melalui \textit{crossover} dan \textit{mutation}, dengan biaya evaluasi yang bergantung pada jumlah populasi dan jumlah generasi. Pada implementasi ini, \textit{crossover} yang menjaga validitas mengeluarkan paket warisan yang bertabrakan atau kehilangan penyangga, sehingga anak cenderung bernilai lebih rendah daripada orang tuanya.

Pengujian dilakukan pada tiga \textit{test case} (\texttt{tc8}, \texttt{tc9}, dan \texttt{tc10}) dengan tiga \textit{run} untuk setiap algoritma atau konfigurasi. Dari sisi kualitas solusi, \textit{Genetic Algorithm} dan \textit{Simulated Annealing} menghasilkan nilai objektif akhir yang lebih tinggi daripada seluruh varian \textit{Hill-climbing}. Nilai terbaik yang ditemukan adalah 103 pada \texttt{tc8}, 114 pada \texttt{tc9}, dan 97 pada \texttt{tc10}. Pada \texttt{tc8}, \textit{Simulated Annealing} dan seluruh konfigurasi \textit{Genetic Algorithm} mencapai 103 pada ketiga \textit{run}, sedangkan rata-rata varian \textit{Hill-climbing} berkisar 97,33 hingga 101,33. Pada \texttt{tc9}, \textit{Genetic Algorithm} dengan populasi 50 atau 100 memperoleh rata-rata 111,33, \textit{Simulated Annealing} memperoleh 105,33, dan varian \textit{Hill-climbing} hanya 99,33 hingga 101,00. Pada \texttt{tc10}, \textit{Genetic Algorithm} dengan populasi 100 konsisten mencapai 97, sedangkan \textit{Simulated Annealing} memperoleh rata-rata 86,00 dan varian \textit{Hill-climbing} 84,00 hingga 89,33. Nilai \textit{global optimum} tidak diketahui secara pasti, sehingga nilai terbaik yang ditemukan dijadikan acuan. Sebagai gambaran, pada \texttt{tc9} total nilai seluruh paket adalah 119 (dan seluruh paket secara total tidak melebihi kapasitas berat), sehingga nilai 114 yang ditemukan mencapai sekitar 95,8\% dari batas atas tersebut.

Dari sisi durasi, \textit{Steepest Ascent} merupakan yang tercepat dengan rata-rata 0,008 hingga 0,016 detik, diikuti \textit{Sideways Move} dan \textit{Random Restart} (0,06 hingga 0,11 detik). \textit{Simulated Annealing} membutuhkan 0,34 hingga 1,62 detik, sedangkan \textit{Genetic Algorithm} membutuhkan 0,10 hingga 0,90 detik bergantung pada jumlah populasi dan iterasi. \textit{Stochastic Hill-climbing} merupakan yang paling lambat (2,30 hingga 12,57 detik) karena pada seluruh \textit{run} menjalankan 10.000 iterasi penuh, tetapi nilai akhirnya tidak lebih baik daripada \textit{Simulated Annealing}. Dari sisi konsistensi, \textit{Genetic Algorithm} dengan populasi besar paling stabil (simpangan baku 0,00 pada \texttt{tc8} dan \texttt{tc10} untuk populasi 100), sedangkan \textit{Random Restart} paling bervariasi pada \texttt{tc10} (simpangan baku 7,09) karena hasilnya bergantung pada titik awal yang terbatas jumlahnya.

Pada \textit{Genetic Algorithm}, jumlah populasi lebih berpengaruh daripada jumlah iterasi. Pada seluruh 45 \textit{run}, individu terbaik selalu berasal dari populasi awal, dan rata-rata \textit{fitness} populasi justru menurun setelah generasi pertama. Akibatnya, dengan populasi 50, menaikkan iterasi dari 50 menjadi 100 dan 200 tidak mengubah nilai akhir pada ketiga \textit{test case} dan hanya menambah durasi. Sebaliknya, populasi 20 menghasilkan nilai yang lebih rendah (rata-rata 102,33 pada \texttt{tc9} dan 90,33 pada \texttt{tc10}) dibandingkan populasi 50 dan 100, karena populasi awal yang kecil memuat lebih sedikit \textit{state} acak sehingga peluang memperoleh \textit{state} bernilai tinggi lebih kecil.

Secara keseluruhan, \textit{Genetic Algorithm} dengan populasi besar unggul dalam kualitas dan konsistensi hasil, terutama berkat populasi awal yang besar dan bukan proses evolusinya, \textit{Simulated Annealing} memberikan keseimbangan yang baik antara kualitas dan durasi, sedangkan varian \textit{Hill-climbing} unggul dalam kecepatan (khususnya \textit{Steepest Ascent}) tetapi mudah terjebak pada \textit{local optimum}, terutama pada \textit{test case} dengan \textit{constraint} yang lebih ketat seperti \texttt{tc9}.

\section{Saran}

Terdapat beberapa hal yang dapat ditingkatkan pada algoritma pencarian. Pertama, \textit{state} awal saat ini dibangkitkan secara acak. Pembangkitan yang lebih terarah, misalnya memasukkan paket berdasarkan rasio nilai terhadap volume atau berat terlebih dahulu, dapat memberikan titik awal yang lebih baik bagi seluruh algoritma. Kedua, operator \textit{neighbor} saat ini hanya terdiri dari \textit{move}, \textit{rotate}, dan \textit{swap}, sehingga mengeluarkan paket dari truk hanya dapat terjadi melalui \textit{swap}. Penambahan operator khusus untuk menambah atau mengeluarkan paket dapat memperluas ruang pencarian yang terjangkau. Ketiga, pada \textit{Hill-climbing with Sideways Move}, penghitung \textit{sideways move} yang kumulatif dapat diganti dengan penghitung \textit{sideways move} berurutan yang di-\textit{reset} setiap kali terjadi perbaikan, agar pencarian tidak berhenti terlalu dini pada permasalahan dengan banyak daerah datar. Keempat, \textit{Stochastic Hill-climbing} selalu menghabiskan seluruh 10.000 iterasi sehingga menjadi algoritma paling lambat; penambahan kriteria berhenti berdasarkan stagnasi nilai objektif dapat mengurangi durasi tanpa mengubah kualitas hasil secara berarti.

Pada \textit{Simulated Annealing}, kualitas hasil sangat dipengaruhi oleh suhu awal, laju pendinginan, dan suhu minimum. Suhu awal sebaiknya disesuaikan dengan skala nilai paket pada setiap \textit{test case}, karena selisih nilai objektif $\Delta E$ menentukan probabilitas penerimaan langkah yang lebih buruk. Skema pendinginan yang adaptif atau mekanisme \textit{reheating} ketika pencarian terjebak pada \textit{local optimum} juga dapat dicoba. Pada \textit{Genetic Algorithm}, dapat dilakukan penyetelan parameter lebih lanjut, terutama probabilitas \textit{crossover} dan \textit{mutation}, serta pengujian metode seleksi lain seperti \textit{tournament selection}. Implementasi saat ini mengganti seluruh populasi pada setiap generasi tanpa \textit{elitism}, sehingga penerapan \textit{elitism} untuk menyimpan individu terbaik pada setiap generasi dapat mencegah hilangnya solusi yang baik. Selain itu, karena \textit{crossover} saat ini membuat anak kehilangan banyak paket, operator \textit{crossover} yang lebih sesuai dengan struktur ruang, misalnya mewarisi satu bagian ruang truk secara utuh dari tiap orang tua, perlu dicoba agar proses evolusi benar-benar dapat menghasilkan individu yang lebih baik daripada populasi awal.

Dari sisi efisiensi, pembangkitan \textit{state} awal dan \textit{neighbor} masih mengenumerasi seluruh koordinat yang mungkin dan memvalidasi seluruh \textit{state} dari awal setiap kali. Pembatasan kandidat posisi pada titik-titik yang bersinggungan dengan dinding atau paket lain, serta validasi inkremental yang hanya memeriksa paket yang berubah, dapat mempercepat pencarian secara signifikan, terutama pada \textit{test case} dengan jumlah paket besar.

Terakhir, pengujian pada tugas ini baru mencakup tiga \textit{test case} berskala kecil (8 hingga 10 paket), sehingga perbedaan kinerja antaralgoritma pada \textit{test case} yang lebih besar belum teruji. Penelitian selanjutnya sebaiknya mengevaluasi algoritma pada lebih banyak \textit{test case} dengan variasi jumlah paket, ukuran truk, dan proporsi paket rapuh, serta menggunakan lebih banyak \textit{run} dengan \textit{seed} berbeda agar konsistensi hasil dapat dinilai lebih andal.